\begin{proof}
If $\Lambda=\mathbf Z$, openness of the semipositivity locus
\cite[Remark~7.3(4)]{GJR26} lets us shrink around $x$ until $h$ is
semipositive on one strict analytic neighbourhood; then the one-term
decomposition $h=1\cdot h$ has all the asserted properties.  We may
therefore suppose that $\Lambda$ is divisible.

By \cite[Proposition~5.8]{GJR26}, after shrinking to a strict analytic
neighbourhood $U$ of $x$ we can write
\[
                         h=\sum_{i=1}^m\lambda_i h_i,
 \qquad \lambda_i\in\Lambda,\quad h_i\in\operatorname{PL}_{\mathbf Z}(U).
\]
Linearity of residues \cite[Lemma~6.18]{GJR26} gives
\[
                         L_h(x)=\sum_{i=1}^m\lambda_iL_{h_i}(x).
\]
The system of models of the reduction of the germ is filtered.  By
\cite[Lemma~6.5 and \S6.6]{GJR26}, we may therefore choose one model $V$ of
$\operatorname{red}(X,x)$ and line bundles $L_i$ on $V$ representing all
the residue classes $L_{h_i}(x)$.  By the hypothesis on $x$,
\[
        \dim V=\operatorname{trdeg}_{\widetilde K}
                    \widetilde{\mathscr H(x)}\leq1.
\]

Suppose first that $\dim V=0$, or that $\dim V=1$ and $V$ is not proper.
Then $V$ contains no integral proper closed curve.  Nefness is consequently
vacuous for every $L_i$ and its dual.  Thus every $h_i$ is harmonic at $x$
by \cite[Proposition~7.1]{GJR26}.  Writing
\[
          \lambda_i=\lambda_i^+-\lambda_i^-,\qquad
          \lambda_i^+=\max\{\lambda_i,0\},\quad
          \lambda_i^-=\max\{-\lambda_i,0\},
\]
expresses $h$ as a finite $\Lambda_{\geq0}$-linear combination of the
integral harmonic functions $h_i$ and $-h_i$.

It remains to treat the case in which $V$ is a proper integral curve.  Put
\[
                 d_i=\deg L_i\in\mathbf Z,\qquad
                 d=(d_1,\ldots,d_m),\qquad
                 D=d\cdot\lambda=\deg L_h(x).
\]
Semipositivity of $h$ says exactly that $D\geq0$.  If $d=0$, every $h_i$
is harmonic at $x$, and the preceding sign-splitting argument applies.
Assume henceforth that $d\ne0$.  Let
\[
                 \delta=\gcd(d_1,\ldots,d_m)>0,
                 \qquad \ell=\delta^{-1}d:\mathbf Z^m\to\mathbf Z.
\]
The homomorphism $\ell$ is primitive.  Choose $u\in\mathbf Z^m$ with
$\ell(u)=1$ and a $\mathbf Z$-basis
$v_1,\ldots,v_{m-1}$ of $\ker\ell$.  Since
$\mathbf Z^m=\mathbf Zu\oplus\ker\ell$, the same integral change of basis
over $\Lambda$ gives a unique decomposition
\[
            \lambda=t u+\sum_{j=1}^{m-1}c_jv_j,
            \qquad
            t=\ell(\lambda)=D/\delta\in\Lambda_{\geq0},
            \quad c_j\in\Lambda.
\]
Define integral PL functions on $U$ by
\[
                    q=\sum_i u_i h_i,
             \qquad r_j=\sum_i(v_j)_i h_i.
\]
Their residue line bundles on $V$ have degrees
\[
              \deg L_q(x)=d\cdot u=\delta>0,
              \qquad \deg L_{r_j}(x)=d\cdot v_j=0.
\]
On a proper integral curve, a line bundle is nef precisely when its degree
is nonnegative and is numerically trivial precisely when its degree is zero.
The residue criterion \cite[Proposition~7.1]{GJR26} therefore shows that
$q$ is semipositive at $x$ and that each $r_j$ and $-r_j$ is harmonic at
$x$.  Now
\[
 h=tq+\sum_j c_j^+r_j+\sum_j c_j^-(-r_j),
 \qquad
 c_j^+=\max\{c_j,0\},\quad c_j^-=\max\{-c_j,0\},
\]
is the required positive-coefficient decomposition at the level of germs.
All summands other than the possibly nonharmonic $q$ are harmonic.

Finally, the semipositivity and harmonicity loci of PL functions are open
by \cite[Remark~7.3(4)]{GJR26}.  Intersecting the finitely many such loci
with $U$ produces one strict analytic neighbourhood $W$ of $x$ on which
every displayed summand has the asserted property.  Restricting the identity
to $W$ completes the proof.
\end{proof}
